%% kapa: exam=1 round=32 year=2021 period=1 form=A subject=economics subject_ko=경제학원론 seq=8 number=48
%% source: Q-Net 감정평가사 기출문제 (cst00309) artlSeq=5208397 file=extracted/1008/5208397/A형/2021년도 제32회 감정평가사 1차 1교시 A형.pdf
%% conversion: kordoc 4.24.51 -> tools/gen1.py (auto); review=proofread-sonnet; flags=none
\begin{question}{48}
기업생산이론에 관한 설명으로 옳은 것을 모두 고른 것은?
\begin{framedtext}
ㄱ. 장기(long-run)에는 모든 생산요소가 가변적이다. \par
ㄴ. 다른 생산요소가 고정인 상태에서 생산요소 투입 증가에 따라 한계생산이 줄어드는 현상이 한계생산 체감의 법칙이다. \par
ㄷ. 등량곡선이 원점에 대해 볼록하면 한계기술대체율 체감의 법칙이 성립한다. \par
ㄹ. 비용극소화는 이윤극대화의 필요충분조건이다. \par
\end{framedtext}
\begin{choices}
  \item ㄱ, ㄴ
  \item ㄷ, ㄹ
  \item ㄱ, ㄴ, ㄷ
  \item ㄴ, ㄷ, ㄹ
  \item ㄱ, ㄴ, ㄷ, ㄹ
\end{choices}
\end{question}
\answer{3}
